% Generated draft: requires mathematical and source review.
\documentclass[11pt]{article}
\usepackage[margin=25mm]{geometry}
\usepackage{fontspec}
\setmainfont{Latin Modern Roman}
\usepackage{amsmath,amssymb,mathtools,unicode-math}
\setmathfont{Latin Modern Math}
\usepackage{enumitem,needspace,xurl,hyperref}
\setlength{\parindent}{0pt}
\setlength{\parskip}{5pt}
\setlength{\emergencystretch}{3em}
\newcommand{\N}{\mathbb{N}}
\newcommand{\Z}{\mathbb{Z}}
\newcommand{\Q}{\mathbb{Q}}
\newcommand{\R}{\mathbb{R}}
\newcommand{\C}{\mathbb{C}}
\newcommand{\F}{\mathbb{F}}
\newcommand{\A}{\mathbb{A}}
\newcommand{\PP}{\mathbb{P}}
\newcommand{\E}{\mathbb{E}}
\newcommand{\eps}{\varepsilon}
\newcommand{\dd}{\,\mathrm{d}}
\begin{document}
\section*{UnsolvedMath ID: 30000955}
\textbf{Source problem code:} catalogue entry
\subsection{Definitions and mathematical statement}
A virtual knot is represented by a one-component knot diagram with classical and virtual crossings, modulo the classical Reidemeister moves and virtual detour moves, with orientation data when specified. Let \(U\) denote the virtual unknot. Write \(B(K)\) for the fundamental biquandle of \(K\), the biquandle presented by semiarc generators and the crossing relations in the chosen biquandle convention. Write \(L(K)\) for the linear-biquandle invariant intended by the report; its coefficient ring, operations, and retained invariant are not specified in this record.
\begin{enumerate}
\item Does there exist a virtual knot \(K\) with \(K\not\simeq U\) that is not detected by \(L\)? If detection means distinguishing from the unknot, this asks for \(L(K)\cong L(U)\) in the intended category of linear invariants. The source proposes investigating the braid closure
\[
K=\widehat{\kappa\tau\kappa^{-1}\tau^{-1}},\qquad \kappa\in\ker(\beta),
\]
where \(\beta\) is the intended Burau map. This is a suggested construction, not an assertion that such a closure is a nontrivial knot or is undetected. The braid group, number of strands, map \(\beta\), and element \(\tau\) are not specified; one must also ensure the closure has a single component.
\item Is the fundamental biquandle complete? For the oriented interpretation in the cleaned statement, the question is whether, for all oriented virtual knots \(K,K'\),
\[
B(K)\cong B(K')\quad\Longrightarrow\quad K\simeq K'
\]
under orientation-preserving virtual-knot equivalence. The original wording instead says ``up to orientation.'' Whether reversal is to be identified, and what orientation or peripheral information is retained in the biquandle isomorphism, must be fixed before this implication is definitive.
\end{enumerate}
\subsection{Short English statement}
Find a nontrivial virtual knot invisible to the intended linear-biquandle invariant, and decide whether the fundamental biquandle determines a virtual knot with the intended orientation convention. The proposed Burau-kernel commutator is a construction to investigate, not a verified example.
\subsection{Sources}
\begin{itemize}
\item Local UnsolvedMath registry, permanent ID \(30000955\), source code \texttt{catalogue entry}, ``Detection and Classification of Virtual Knots by Biquandles.'' Registry curation: Ulam AI / UnsolvedMath Contributors, CC BY 4.0; underlying sources retain their supplied provenance.
\item \emph{Invariants in Low-Dimensional Topology} (2009), Oberwolfach Reports, pp.\ 1157--1232, DOI \url{https://doi.org/10.4171/owr/2008/22}.
\item Supplied ``Hrencecin--Kauffman 2007'': \url{https://arxiv.org/abs/math/0703216}.
\item Supplied ``biquandle status survey'': \url{https://www.sciencedirect.com/science/article/pii/S0022404921002772}.
\item The supplied literature assessment, checked 2026-08-21, labels the item partially solved: braid-level nonfaithfulness phenomena for linear invariants are known, while the requested knot example and completeness question remain unresolved according to that assessment. Braid-level nonfaithfulness alone is not asserted to produce the requested knot.
\end{itemize}
\end{document}
